\chapter{Project Overview and Research Claim}

\begin{introduction}
\item Explain the physical architecture before presenting the optimization.
\item Separate the mathematical relaxation from the certified physical solution.
\item Use the final evidence package as the sole source of numerical claims.
\end{introduction}

\section{Problem in one paragraph}
This project studies downlink resource allocation for a participation-constrained
Cell-Free integrated sensing and communication (ISAC) network. Distributed APs
jointly transmit communication streams to single-antenna UEs, while dedicated
target-specific sensing waveforms are emitted only by selected AP--target
pairs. The central processor minimizes total transmit power subject to robust
communication QoS, sensing-detection QoS, PCRB tracking accuracy, per-AP power,
and sensing-cluster cardinality constraints. Imperfect CSI and binary AP--target
participation make the design a nonconvex mixed-integer problem.

\section{Core contribution}
The contribution is not a global MINLP solver. It is a physically interpretable
and certifiable solution pipeline: covariance lifting and SDR expose an SDP
structure; robust SINR and PCRB constraints are represented by LMIs; double DC
penalties encourage rank-one and binary structure; and a separate recovery
stage fixes a binary topology, re-optimizes continuous covariances, and
validates every physical constraint.

\section{One-page oral summary}
\quick{We propose a Cell-Free ISAC design in which communication remains
globally cooperative, while dedicated sensing waveforms are authorized by a
target-centric AP cluster. The resulting robust mixed-integer design is handled
by SDR, the S-procedure, PCRB Schur LMIs, and double DC-SCA. A fixed-assignment
recovery and validation stage converts the relaxed result into a physically
feasible binary sensing solution.}
\logic{先讲架构的非对称性：通信全网协作，感知按目标选择集群。再讲难点：
鲁棒约束、PCRB 和二元关联耦合。最后强调贡献是\textbf{可认证的启发式恢复流程}，
不是声称原始混合整数问题的全局最优。}

\section{Claim boundaries}
\begin{itemize}
\item SDR is a continuous lower bound, not a deployable design.
\item A solver failure or time limit is not evidence of physical infeasibility.
\item Fig.~8 is a QoS-requirement sensitivity study, not a global Pareto frontier.
\item M12 results demonstrate scalability evidence and computational cost, not real-time operation.
\end{itemize}

\sourcebox{Primary sources: \texttt{PROBLEM\_FORMULATION.md},
\texttt{MASTER\_CONSTRAINTS.md}, and
\texttt{deliverables/paper\_results\_and\_figures\_v1\_0/README.md}.}
